\section{Round-twenty-nine reconstruction: loss-indexed current jets and continuous Hilbert limits}
\label{sec:r29-a1}

This revision replaces the incorrectly directed Sobolev-dual scale of the
previous manuscript.  It also distinguishes shell boundary currents from
cumulative boundary currents and uses a genuinely continuous interpolation
in the functional limit theorem.  No density between mutually singular
infinite Bernoulli products is introduced.

\subsection{The loss-indexed current scale}

Let $\widehat S_a$ be the labelled resolution of the four-branch symplectic
baker section and let $\mathscr F_n$ be the finite collection of bulk, face,
corner, and flag strata created by the first $n$ seam preimages.  Parent-fold
coordinates identify these collections for $a$ in a compact parameter
interval.  For $s\ge s_0$ define
\[
 \Phi^s=\left\{\phi=(\phi_F):
   \sum_{n\ge0}\beta^{-n}\sum_{F\in\mathscr F_n}
   \norm{\phi_F}_{H^{s+\operatorname{codim}F+2}(F)}^2<\infty\right\},
 \qquad \mathbb J^s=(\widehat\Phi^s)',                    \tag{A1.1}
\]
where $\widehat\Phi^s$ is the closed kernel of all trace, incidence, and
subdivision relations.  If $s'>s$, then
$\widehat\Phi^{s'}\hookrightarrow\widehat\Phi^s$ and hence
$\mathbb J^s\hookrightarrow\mathbb J^{s'}$.  Thus increasing the index on
$\mathbb J^s$ means allowing a more singular current.

Let $\nabla_a^{\rm test}$ denote material differentiation in the transported
parent-fold chart.  Its coefficients contain at most two spatial derivatives,
so
\[
 \nabla_a^{\rm test}:\widehat\Phi^{s+2}\longrightarrow\widehat\Phi^s
 \quad\hbox{boundedly}.                                  \tag{A1.2}
\]
The current connection is defined by transposition,
\[
 \ip{\nabla_a J}{\phi}
 =\partial_a\ip{J}{\phi}-\ip{J}{\nabla_a^{\rm test}\phi},
 \qquad \phi\in\widehat\Phi^{s+2}.                       \tag{A1.3}
\]

\begin{lemma}[Correct direction of distributional loss]
\label{lem:r29-a1-direction}
For every $s\ge s_0$,
\[
 \nabla_a:\mathbb J^s\longrightarrow\mathbb J^{s+2}
 \quad\hbox{is closed and bounded on transported compact charts}.       \tag{A1.4}
\]
Consequently the $j$th parameter derivative of a current initially in
$\mathbb J^m$ is interpreted in $\mathbb J^{m+2j}$.
\end{lemma}

\begin{proof}
For $J\in\mathbb J^s$ and $\phi\in\widehat\Phi^{s+2}$, (A1.2) gives
\[
 \abs{\ip{J}{\nabla_a^{\rm test}\phi}}
 \le C\norm{J}_{\mathbb J^s}\norm{\phi}_{\widehat\Phi^{s+2}}.
\]
The ordinary derivative of the transported coefficient functional has the
same estimate.  Formula (A1.3) therefore defines an element of
$(\widehat\Phi^{s+2})'=\mathbb J^{s+2}$.  Closedness follows by passing to the
limit in (A1.3).  This direction agrees with the model
$\delta_x\mapsto\partial_x\delta_x$: differentiation requires a larger, not a
smaller, negative Sobolev space.
\end{proof}

\subsection{Coherent shell and cumulative boundary data}

Let $\Sigma=\{1,2,3,4\}^{\mathbb Z}$ and
$\mu_a=\bigotimes_{k\in\mathbb Z}w(a)$, with all symbol weights positive and
$C^{r+2}$.  Write $\mathcal F_N=\sigma(\omega_k:|k|\le N)$ and
$\Pi_N^a=E_{\mu_a}[\cdot\mid\mathcal F_N]$.

A coherent current jet of order $r$ consists of
\[
 J^{(j)}_{a,N}\in L^2(\mu_a|_{\mathcal F_N};\mathbb J^{m+2j}),
 \qquad 0\le j\le r,                                     \tag{A1.5}
\]
with $\Pi_N^aJ^{(j)}_{a,N+1}=J^{(j)}_{a,N}$.  The geometric derivative at
level $N$ is decomposed as
\[
 \partial_aJ^{(j)}_{a,N}=J^{(j+1)}_{a,N}+B^{(j)}_{a,N},
 \qquad B^{(j)}_{a,N}=\sum_{q=0}^{N}b^{(j)}_{a,q}.         \tag{A1.6}
\]
Here $b^{(j)}_{a,q}$ is the new boundary shell born at exact depth $q$ and
$B^{(j)}_{a,N}$ is cumulative.  Only the shells are required to decay.

\begin{definition}[Loss-indexed projective response domain]
\label{def:r29-a1-domain}
For $\eta>0$, let $\mathfrak C^{r}_{\eta}$ be the coherent jets satisfying
\[
 \sum_{j=0}^{r}\sup_a\left[
   \sup_N\norm{J^{(j)}_{a,N}}_2+
   \sum_{N\ge0}e^{\eta N}
      \norm{J^{(j)}_{a,N+1}-J^{(j)}_{a,N}}_2+
   \sum_{q\ge0}e^{\eta q}\norm{b^{(j)}_{a,q}}_2
 \right]<\infty.                                         \tag{A1.7}
\]
The normally convergent cumulative boundary current is
$B^{(j)}_a=\sum_{q\ge0}b^{(j)}_{a,q}$.
\end{definition}

For $M>N$ put
\[
 S_{a;N,M}=\sum_{N<|k|\le M}\partial_a\log w_{\omega_k}(a).
\]

\begin{theorem}[Projective derivative with coherent boundary shells]
\label{thm:r29-a1-projective}
For $\mathbf J\in\mathfrak C^r_\eta$, the projective limit
$J^{(j)}_a=\lim_NJ^{(j)}_{a,N}$ exists in
$L^2(\mu_a;\mathbb J^{m+2j})$.  Moreover
\begin{equation}\tag{A1.8}
\begin{split}
 \partial_a\Pi_N^aJ^{(j)}_a
  ={}&\Pi_N^aJ^{(j+1)}_a+\Pi_N^aB^{(j)}_a\\
    &+\lim_{M\to\infty}\Pi_N^a\!\bigl[
      (J^{(j)}_{a,M}-\Pi_N^aJ^{(j)}_{a,M})S_{a;N,M}\bigr]\\
    &+\Pi_N^a\!\bigl[(J^{(j)}_a-\Pi_N^aJ^{(j)}_a)
      S_{a;-N,N}\bigr],
\end{split}
\end{equation}
where every term lies in $\mathbb J^{m+2j+2}$.  Formula (A1.8) may be
differentiated through order $r-j-1$, with normal convergence in the
corresponding loss-indexed current space.
\end{theorem}

\begin{proof}
The first two series in (A1.7) make $(J^{(j)}_{a,N})_N$ a Cauchy martingale and
make $B^{(j)}_a$ well defined.  On a shell of width $q$, the centred product
score has $L^2$ norm $O(q^{1/2})$, while the current martingale difference is
$O(e^{-\eta q})$.  Cauchy--Schwarz gives the summable majorant
$C(1+q)^{1/2}e^{-\eta q}$.  Higher derivatives introduce only fixed-degree
Bell polynomials in the shell score, hence polynomial factors absorbed by the
same exponential series.  Finite-cylinder differentiation followed by
$M\to\infty$ proves (A1.8).  Lemma~\ref{lem:r29-a1-direction} determines the
correct target space after each derivative.
\end{proof}

\subsection{A non-formal geometric certificate for the affine benchmark}

For a word $\omega_0\ldots\omega_{n-1}$ let
$p_\omega(a)=\prod_{k<n}w_{\omega_k}(a)$.  A depth-$n$ seam sheet has a
uniformly bounded affine incidence multiplicity $d_0$ and its current norm in
$\mathbb J^{m+2j}$ is bounded by $C_jn^jp_\omega(a)$.  Therefore
\[
 \sum_{|\omega|=n}
 \norm{\nabla_a^j[J_a^{\rm seam}]_\omega}_{\mathbb J^{m+2j}}^2
 \le C_jn^{2j}\sum_{|\omega|=n}p_\omega(a)^2
 =C_jn^{2j}\left(\sum_{i=1}^{4}w_i(a)^2\right)^n.         \tag{A1.9}
\]
The same estimate holds for every trace-incidence image, with an additional
factor bounded by $d_0^2$.

\begin{proposition}[Affine seam-current membership]
\label{prop:r29-a1-geometric}
Assume
\[
 \sup_a\sum_{i=1}^{4}w_i(a)^2=\vartheta^2<1
 \quad\hbox{and}\quad \beta d_0^2<\vartheta^{-2}.         \tag{A1.10}
\]
Then the labelled seam current and all of its material derivatives through
order $r$ belong to $\mathfrak C^r_\eta$ for every
$0<\eta<-\log\vartheta$.  The boundary shell at depth $q$ satisfies
\[
 \norm{b^{(j)}_{a,q}}_2\le C_j(1+q)^j\vartheta^q.          \tag{A1.11}
\]
\end{proposition}

\begin{proof}
Equation (A1.9) is an exact product computation, not an appeal to a branch
width slogan.  The factor $d_0$ controls every incidence row and (A1.10)
absorbs the depth weight in (A1.1).  The difference between cylinder depths
$q-1$ and $q$ consists exactly of the newly born depth-$q$ sheets, so (A1.11)
is a shell estimate.  Polynomial factors from parameter derivatives remain
summable against $e^{\eta q}\vartheta^q$.
\end{proof}

\subsection{Continuous interpolation and differentiable covariance}

Let
$Y_k=J^{(0)}_a\circ\sigma^k-EJ^{(0)}_a$.  For the past filtration define
$P_0Y_q=E[Y_q\mid\mathcal F_0^-]-E[Y_q\mid\mathcal F_{-1}^-]$ and suppose
\[
 \sum_{q\in\mathbb Z}(1+|q|)^{1/2}
\norm{P_0Y_q}_{L^2(\mathbb J^m)}<\infty.
                                                                    \tag{A1.12}
\]
Set $D_0=\sum_qP_0Y_q$ and $D_k=D_0\circ\sigma^k$.

\begin{lemma}[Uniform maximal martingale approximation]
\label{lem:r29-a1-maximal}
Under (A1.12),
\[
 n^{-1/2}\max_{j\le n}
 \norm{\sum_{k<j}(Y_k-D_k)}_{\mathbb J^m}\longrightarrow0
 \quad\hbox{in }L^2,                                      \tag{A1.13}
\]
and
\[
 n^{-1/2}\max_{k<n}\norm{Y_k}_{\mathbb J^m}
 \longrightarrow0\quad\hbox{in probability}.             \tag{A1.14}
\]
\end{lemma}

\begin{proof}
The dyadic projection proof of the Maxwell--Woodroofe estimate gives (A1.13):
truncate $D_0$ at range $2^L$, write the finite-range remainder as a
coboundary, and control the projection tail by Doob's inequality and the
weighted series (A1.12).  For (A1.14), uniform integrability of
$\norm{Y_0}^2$ gives
\[
 P\{\max_{k<n}\norm{Y_k}>\epsilon\sqrt n\}
 \le nP\{\norm{Y_0}>\epsilon\sqrt n\}=o(1).
\]
\end{proof}

Define the actual polygonal interpolation
\[
 \widetilde W_n(t)=n^{-1/2}\left[
  \sum_{k<\lfloor nt\rfloor}Y_k+
  (nt-\lfloor nt\rfloor)Y_{\lfloor nt\rfloor}\right].     \tag{A1.15}
\]

\begin{theorem}[Trace-class current FCLT with response]
\label{thm:r29-a1-fclt}
The process $\widetilde W_n$ is a random element of
$C([0,1];\mathbb J^m)$ and converges there to a Brownian motion with covariance
operator
\[
 Q_a h=E[\ip{D_0}{h}D_0],\qquad
 \operatorname{tr}Q_a=E\norm{D_0}_{\mathbb J^m}^2.         \tag{A1.16}
\]
If the projective jet belongs to $\mathfrak C^r_\eta$, then
$a\mapsto Q_a$ is $C^{r-1}$ in trace norm.  Its derivatives are obtained by
termwise differentiation of the normally convergent projection series for
$D_0$.
\end{theorem}

\begin{proof}
The Hilbert-valued martingale FCLT applies to $(D_k)$; finite-dimensional
conditional brackets converge by ergodicity, and the trace identity gives
tightness of Hilbert tails.  Lemma~\ref{lem:r29-a1-maximal} transfers the
limit to the partial sums.  The sup norm between the step interpolation and
(A1.15) is bounded by the left side of (A1.14), so the continuous path-space
claim follows.

For covariance response, differentiate each finite projection
$P_0Y_q$ using Theorem~\ref{thm:r29-a1-projective}.  The derivatives are
bounded by $C(1+|q|)^re^{-\eta|q|}$ in the appropriate current space.
Products entering $D_0\otimes D_0$ are therefore normally summable in trace
norm.  Dominated differentiation proves the last assertion.
\end{proof}

The repaired construction has the correct distributional direction, coherent
boundary data, an explicit affine membership estimate, a genuinely continuous
interpolation, and a trace-norm differentiable covariance family.
